\mxnewpage\section*{Correction EV1-B}\label{EV1-B-CORR}
\textbf{QCM :} 1 B, 2 B, 3 D, 4 A, 5 D, 6 C, 7 C, 8 B, 9 A, 10 C, 11 B, 12 C. Chaque réponse correcte : 0{,}5 point ; total 6.
\paragraph{Vérifications des automatismes.} Priorités ; somme de fractions ; soustraction des exposants ; coefficient multiplicateur 1{,}2 ; division par 0{,}8 ; produit 1{,}25 fois 0{,}8 ; résolution affine ; accroissement vertical divisé par l’horizontal ; moyenne ; angle divisé par 360 ; complément à 1 ; substitution entre parenthèses.
\subsection*{Exercice 1 -- Une croissance par motif}
\begin{enumerate}
\item $3,8,15$. \textit{Barème : 1 point. Résultat et justification élémentaire attendus ; demi-point possible pour démarche correcte inachevée.}
\item Différence $2n+3>0$ : stricte croissance. \textit{Barème : 2 points. Méthode ou justification : 1 ; calcul et conclusion : 1.}
\item $n^2+2n-99=(n-9)(n+11)\geq0$ avec $n\geq1$ : $n\geq9$ ; $u_9=99$. \textit{Barème : 2 points. Méthode ou justification : 1 ; calcul et conclusion : 1.}
\end{enumerate}
\subsection*{Exercice 2 -- Deux dimensions et une contrainte}
\begin{enumerate}
\item $x+y=13$, $xy=40$, x,y positifs. \textit{Barème : 1 point. Résultat et justification élémentaire attendus ; demi-point possible pour démarche correcte inachevée.}
\item Racines de $t^2-13t+40=(t-5)(t-8)$ : 5 m et 8 m. \textit{Barème : 2 points. Méthode ou justification : 1 ; calcul et conclusion : 1.}
\item $A(x)=x(13-x)=\dfrac{169}{4}-(x-\dfrac{13}{2})^2$, maximum $\dfrac{169}{4}=42{,}25$ m$^2$. Écart $\dfrac{9}{4}$ m$^2$. \textit{Barème : 2 points. Méthode ou justification : 1 ; calcul et conclusion : 1.}
\end{enumerate}
\subsection*{Exercice 3 -- Angle puis approximation}
\begin{enumerate}
\item $\cos \theta =\dfrac{3}{6}=\dfrac{1}{2}$, donc $60^\circ$. \textit{Barème : 1 point. Résultat et justification élémentaire attendus ; demi-point possible pour démarche correcte inachevée.}
\item Carré $4+9-6=7$ ; norme $\sqrt7$. \textit{Barème : 1 point. Résultat et justification élémentaire attendus ; demi-point possible pour démarche correcte inachevée.}
\item Approximation 1{,}06 ; exact 1{,}0609 ; erreur 0{,}0009. Cette question indépendante vérifie la dérivation locale. \textit{Barème : 2 points. Méthode ou justification : 1 ; calcul et conclusion : 1.}
\end{enumerate}
\paragraph{Étalonnage.} Durée cible : 20–25 min pour les automatismes, 80–85 min pour les exercices, 10–15 min de relecture. Ces durées sont des estimations éditoriales à mesurer en classe. Accepter toute méthode correcte ; ne pas sanctionner deux fois une erreur numérique propagée avec cohérence. Le barème proposé ne vaut pas décision d’harmonisation de l’équipe.
